% Solution: Gibbons (1992) Problem 3.4 (two games drawn by Nature). Input by \solitem.
% Rewritten from ../pratica-lecture-07-solutions/sol-gibbons-3-4.tex, deepened. Check: check_F_L06.py.
\begin{exsol}{Gibbons Problem 3.4}{169--170}{Exam level}{2}
\asked All pure-strategy BNE when Nature picks Game~1 or Game~2 with probability $\frac12$ each,
player~1 learns which game was drawn, player~2 does not, and they move simultaneously.

\method The recipe of Lecture~6, frame~36: (1) each type of the informed player best-responds;
(2) the uninformed player maximizes expected payoff with his belief; (3) keep profiles in which
everyone best-responds at once. Strategies are type-contingent (frame~29): player~1's strategy is a
pair $s_1 = (s_1(G_1), s_1(G_2)) \in \{T,B\}^2$, written $TT, TB, BT, BB$. Player~2 has one type.
Book: Gibbons \S3.1.B--C (pp.~146--152).

\solution
\textbf{Step 1: player~1, type by type.}
\begin{center}
\begin{tabular}{l|cc}
& against $L$ & against $R$\\\hline
Game~1 type & $T$ ($1 > 0$) & $T$ or $B$ ($0 = 0$)\\
Game~2 type & $T$ or $B$ ($0 = 0$) & $B$ ($2 > 0$)
\end{tabular}
\end{center}
So $BR_1(L) = \{TT, TB\}$ and $BR_1(R) = \{TB, BB\}$. The strategy $BT$ is never a best reply.

\textbf{Step 2: player~2.} He gets a positive payoff only in two cells: $(T,L)$ in Game~1 (payoff
$1$) and $(B,R)$ in Game~2 (payoff $2$). With probability $\frac12$ on each game,
\[
  EU_2(L \mid s_1) = \tfrac12\cdot\mathbf 1[s_1(G_1) = T], \qquad
  EU_2(R \mid s_1) = \tfrac12\cdot 2\cdot\mathbf 1[s_1(G_2) = B].
\]
\begin{center}
\begin{tabular}{l|cccc}
$s_1$ & $TT$ & $TB$ & $BT$ & $BB$\\\hline
$EU_2(L)$ & $\frac12$ & $\frac12$ & $0$ & $0$\\
$EU_2(R)$ & $0$ & $1$ & $0$ & $1$\\
$BR_2$ & $L$ & $R$ & $L, R$ & $R$
\end{tabular}
\end{center}

\textbf{Step 3: mutual best replies.}
\begin{itemize}
  \item $(TT, L)$: $TT \in BR_1(L)$ and $BR_2(TT) = L$. \checkmark
  \item $(TB, L)$: $TB \in BR_1(L)$ but $BR_2(TB) = R$. \texttimes
  \item $(TB, R)$: $TB \in BR_1(R)$ and $BR_2(TB) = R$. \checkmark
  \item $(BB, R)$: $BB \in BR_1(R)$ and $BR_2(BB) = R$. \checkmark
  \item $BT$ is in neither $BR_1(L)$ nor $BR_1(R)$, so no profile with $BT$ is an equilibrium.
\end{itemize}
Equivalently, in the $4 \times 2$ Bayesian normal form (ex-ante payoffs, each game weighted
$\frac12$):
\begin{center}
\begin{tabular}{l|cc}
& $L$ & $R$\\\hline
$TT$ & $\br{\tfrac12}, \br{\tfrac12}$ & $0, 0$\\
$TB$ & $\br{\tfrac12}, \tfrac12$ & $\br{1}, \br{1}$\\
$BT$ & $0, \br{0}$ & $0, \br{0}$\\
$BB$ & $0, 0$ & $\br{1}, \br{1}$
\end{tabular}
\end{center}
The cells with both payoffs underlined are $(TT,L)$, $(TB,R)$ and $(BB,R)$.

\interp Both players want to match the game: $(T,L)$ in Game~1, $(B,R)$ in Game~2. Player~1 knows
the game but player~2 does not, and player~2 must commit to one column. The three equilibria are
three ways to resolve this coordination problem: player~2 bets on Game~1 ($TT, L$) or on Game~2
($TB, R$ and $BB, R$). In $(TB, R)$ and $(BB, R)$ the Game~1 type is indifferent, which is why both
of its actions survive. $(TB,R)$ and $(BB,R)$ give player~2 an expected payoff of $1$, more than
$\frac12$ in $(TT,L)$: the high-stakes Game~2 is worth coordinating on.

\verify \texttt{check\_F\_L06.py} builds the $4\times 2$ Bayesian normal form from the two matrices
and enumerates all eight profiles; it also checks the type-by-type test. Output: ``Gibbons 3.4: pure
BNE = (TT,L), (TB,R), (BB,R)''. Gibbons gives no answer.

\mistakes
\begin{itemize}
  \item Dropping $(TB,R)$ or $(BB,R)$ because the Game~1 type is ``not doing anything''. Ties are best
  replies: when player~2 plays $R$ the Game~1 type gets $0$ either way, so both $T$ and $B$ qualify.
  \item Solving each game separately and reporting $(T,L)$ and $(B,R)$ as ``the equilibria''. Player~2
  cannot play $L$ in Game~1 and $R$ in Game~2: he does not know which game he is in.
  \item Computing player~2's payoff from the realized game only. He must average over both games
  with probability $\frac12$, using player~1's whole strategy.
  \item Writing player~1's strategy as one action: in a Bayesian game it is one action per type.
\end{itemize}

\answer{Three pure-strategy BNE (player~1's strategy as action in Game~1, action in Game~2):
$(TT, L)$, $(TB, R)$ and $(BB, R)$.}
\end{exsol}
